\documentclass[../main.tex]{subfiles}
\begin{document}
%==========================================TIÊU ĐỀ TẬP========================================
\begin{table}[h]
    \begin{adjustwidth}{-1cm}{}
        \begin{tabular}{>{\centering\arraybackslash}p{6cm}|>{\raggedright\arraybackslash}p{12.5cm}}
            \multirow{5}{6cm}
            % Cột bên trái: ÉPISODE và số tập
            {\rainbowcolor
                {\\[-40pt]
                    \flaregothic\fontsize{20pt}{20pt}\selectfont \centering
                    \color{eptype}JOUR\color{black}\\[12pt]
                    \hbrk\fontsize{72pt}{72pt}\selectfont
                    \color{epnum}5
                    \\[18pt]
                    % \flaregothic\fontsize{14pt}{14pt}\selectfont
                    % \color{epattr}BIENVENUE, \uline{\color{Banpire} Mel} !
                    % \flaregothic\fontsize{18pt}{18pt}\selectfont
                    % \color{epattr}ÉPISODE PERDU\\
                    \unrainbowcolor}
                \color{black}}
             &
            % Dòng thứ nhất: tên tập tiếng Nhật
            {\fotskip\fontsize{18pt}{18pt}\selectfont \ruby{東}{とう}\ruby{京}{きょう}の\ruby{夜}{よ}\ruby{空}{ぞら}を\ruby{彩}{いろど}る…
            }                     \\[6pt]
            % Dòng thứ hai: tên tập tiếng Anh
             & {\cafeteria\fontsize{18pt}{18pt}\selectfont How NOT to Behave in the Library}                                               \\[4pt]
            % Dòng thứ ba: tên tập tiếng Việt
             & {\vnmsans\fontsize{18pt}{18pt}\selectfont Bài thơ lục bát về những trò đùa trong thư viện}                                  \\[8pt]
            % Dòng thứ tư: Nhân vật xuất hiện
             & {\sen{\fontsize{14pt}{14pt}\selectfont Track used: Earth Core Adventure (\ruby{地}{ち}\ruby{心}{しん}\ruby{探}{たん}\ruby{検}{けん})}} \\[6pt]
            % Dòng cuối cùng: Thời gian diễn ra của tập (thường sẽ là ngày phát hành)
             & {\continuum\fontsize{16pt}{16pt}\selectfont Ngày phát hành: 03/07/2021}                                                     \\[18pt]
        \end{tabular}
    \end{adjustwidth}
\end{table}
%=========================================TRUYỆN CHÍNH========================================
\
\begin{wrapfigure}[28]{l}{5cm}
    \includegraphics[width=5cm]{../../../../Image/Book?/banderas_model.png}
\end{wrapfigure}

\color{black}

\begin{center}(\uline{Do tập này được viết dạng bài thơ, nên phần giới thiệu nhân vật xuất hiện trong tập này sẽ đặt đầu tiên, để phần kể chuyện sẽ chỉ tập trung vào phần thơ.})\end{center}

\chrint

Belmond Banderas (ベルモンド・バンデラス, Be-rư-mơn Ban-đê-rát(-xừ)) là một thành viên thuộc thế hệ thứ Hai của Nijisanji SEEDs. Một người đã tồn tại từ khi cú nổ Big Bang xảy ra, tức là đâu đó từ 4,6 tới 13,8 tỷ năm về trước, hay theo lời của ông ấy, ``đã  tồn tại từ khi sinh vật sống bắt đầu nhìn thấy những giấc mơ.'' Một ông chú có tính cách thẳng tính và đôi khi cũng trở nên nóng tính.

Ông ấy là chủ quán bar mang tên ``Bar Deras'' (đặt theo tên của mình), nằm ở một nơi mà bất cứ ai đang mơ đều có thể ghé qua. Quán thường tiếp đón khách thuộc nhiều chủng tộc khác nhau tùy theo sở thích của ông ấy, vì vậy đây cũng là địa điểm quen thuộc của các thành viên nhà Cầu Vồng. Banderas được mọi người và người hâm mộ yêu mến nhờ cách trò chuyện đầy cuốn hút.

Ông nổi tiếng là người ga-lăng (theo lời kể của đồng nghiệp), từng đãi ăn một nhóm người trong công ty khi họ trốn trong nhà vệ sinh - sau này mới biết số tiền đó là do người hâm mộ quyên góp chứ không phải của ông ấy.

Banderas có giọng hát truyền cảm và thường phát trực tiếp từ nửa đêm đến sáng, bởi ông ấy là bartender xuất hiện trong giấc mơ. Nhờ giọng hát và nhạc nền jazz êm dịu, người hâm mộ thường nghe stream của ông ấy để dễ ngủ. Ông ấy luôn cảm thấy vui khi có thể giúp đỡ mọi người bằng bất cứ cách nào.

Dù vậy, ông ấy cũng có một mặt khác, đôi khi trở thành một kẻ man rợ hoặc một ông trùm mafia.

Ngoài ra, mọi người ở Nijisanji thường đùa rằng ông ấy là một ``ông chú não tàn'' hoặc ``người chỉ biết vâng dạ'' vì ông ấy luôn đồng ý với mọi lời người khác nói. Đôi khi ông ấy còn hành động như một nữ sinh trung học, với giọng nói cao vút và cách nói chuyện như một nữ sinh khi ông ấy phấn khích.

\begin{wrapfigure}[31]{r}{7.5cm}
    \includegraphics[width=7.5cm]{../../../../Image/Book?/kai_model.png}
\end{wrapfigure}

Belmond có mái tóc đuôi ngựa màu nâu sẫm buộc lệch sang bên trái. Ông có đôi mắt vàng và chùm râu lông nâu ở dưới cằm. Bộ trang phục của ông là bộ đồ pha chế trong các quán rượu thường thấy ở các trò chơi nhập vaai: quần đen, áo da nâu có tay áo màu đen và cổ tay áo xanh nước biển, phía trước mặc tạp dề đen buộc từ hông trở xuống, đi giày da màu nâu. Belmond đeo một chiếc vòng cổ bạc có họa tiết ngôi sao, một chiếc vòng tay nâu với mặt dây chuyền hình ngôi sao, và một sợi dây thắt lưng da nâu với mặt dây chuyền hình ngôi sao.

Mayuzumi Kai (\ruby{黛}{まゆずみ}\ruby{灰}{かい}, \ruby{Ma-giu-dư-ki}{Đại}\ \ruby{Khôi}{Cai}) là một thành viên khác gia nhập vào Nijisanji năm 2019. Cậu chàng 22 tuổi này là một hacker mũ trắng gia nhập vào công ty này với mục đích làm nhiệm vụ mật. Cậu là một người hiếm khi ra ngoài, tới độ khi mua hàng thì luôn mua trực tuyến tất cả mọi thứ, kể cả là nhu yếu phẩm.

Kai là một người điềm tĩnh và kín đáo, đôi khi còn bị gắn mác là ``trầm lặng'' và ``ít nói.'' Vậy nên, bất cứ khi nào cậu ta cười vì một vấn đề nào đó, tất cả mọi người đều\dots\ phấn khích tới lạ thường. Tuy nhiên, Kai là một người rất hòa nhã và nói chuyện tốt với những người cùng công ty, bao gồm cả những người ngoài Nijisanji như Matsuri (cô nàng Lễ hội hướng ngoại).

Cậu chàng nổi tiếng với lối nói chuyện đều đều và có khiếu hài hước theo kiểu thẳng thừng mà mỉa mai. Cách xưng hô của cậu đôi lúc bị nhận xét là ``trống không'', khiến họ rất khó chịu nhưng chẳng làm gì được. Cậu thường đột nhiên hành động khác thường thông qua việc lồng tiếng cường điệu hoặc chỉnh sửa video kỳ quặc; nổi tiếng nhất là, theo gợi ý của khán giả, việc đã hoàn toàn đảo ngược tính cách của mình và hành động một cách năng nổ và lịch sự khiến cả công ty một phen náo loạn.

Mayuzumi Kai đã rời Nijisanji vào 28 tháng Bảy năm 2022, với lý do là không cùng định hướng phát triển của công ty. Bài đăng cuối cùng trên Twitter (nay là X) có nội dung giống như những người thông thạo bảo mật thông tin sẽ hiểu:

\txtbx{TH3 3ND. TH3 2TORY DO32 NOT CONT1NU3.\\(Hết. Câu chuyện không còn được tiếp tục.)}

Kai có mái tóc đen ngắn, ở bên trái có nhuộm một vài đường xanh lục nhạt, đeo khuyên tai đen và đôi mắt xanh lục nhạt. Do luôn ở nhà, Kai thường thấy trong bộ trang phục với quần bó màu xám, áo phông đen giản đơn và đi dép lê trắng hình con thú. Khi ra ngoài, cậu có khoác thêm một chiếc áo màu trắng cổ tay màu đen, viền xanh, thường mặc trễ vai xuống.

\begin{wrapfigure}[30]{l}{9cm}
    \includegraphics[width=9cm]{../../../../Image/Book?/akina_model.jpg}
\end{wrapfigure}

Saegusa Akina (\ruby{三}{さえ}\ruby{枝}{ぐさ}\ruby{明}{あき}\ruby{那}{な}, \ruby{Sa-ê}{Tam}-\ruby{gư-xa}{Chi}\ \ruby{A-ki}{Minh}-\ruby{na}{Na}) là một chàng trai khác gia nhập Nijisanji cùng thời điểm với Kai. Một chàng sinh viên ngoài hai mươi đầy nhiệt huyết với hoài bão lớn và tận tâm với mọi việc mình làm. Cậu chàng này mơ ước lớn nhất là được biểu diễn trực tiếp tại Budoukan.

Cậu có tính cách thân thiện và hoạt bát. Akina nổi tiếng là người ồn ào và dễ bị kích động, dễ bị cuốn vào bất kỳ trò chơi nào mà cậu đang chơi. Đúng với tiểu sử của mình, cậu là một ca sĩ rất năng động: cậu thường tham gia cover các bài hát của các thành viên khác, tổ chức các buổi phát trực tiếp karaoke và tổ chức chương trình Nijimen Utaite Relay đầu tiên. Tuy thế, cậu bạn này lại\dots\ không giỏi ở khoản chơi các nhạc cụ.

Cậu thường bị mọi người ở Nijisanji trêu là ``một chàng trai còn trinh'' và cũng tự nhận mình là một shotacon. Trước những người không quen biết, cậu dễ dàng bối rối - đặc biệt là với phụ nữ. Mặc dù đã quen biết nhiều phụ nữ (chủ yếu là ở cùng công ty) được 8 tháng, nhưng trong lần hợp tác ngoại tuyến đầu tiên với Aizono Manami, anh đã rất khó khăn để nhìn thẳng vào cô ấy - và đó cũng là lý do vì sao cậu thường trốn ở các buổi collab trực tiếp.

Akina có mái tóc đen với chỏm tóc ngố và nhuộm vài sợi đỏ ở bên trái, cùng phong cách như Kai, cùng với đôi mắt xanh nước biển, tựa giống như Kaigou. Cậu chàng này khi ra ngoài (bất kể là đi học hay đến công ty) thường mặc áo trùm đầu màu trắng bên trong với áo khoác đen bên ngoài, đi cùng với quần đen và giày thể thao trắng.

\begin{wrapfigure}[27]{r}{12cm}
    \includegraphics[width=12cm]{../../../../Image/Book?/ritsuki_model.jpg}
\end{wrapfigure}

Sakura Ritsuki (\ruby{桜}{さくら}\ruby{凛}{り}\ruby{月}{つき}, \ruby{Xa-cư-ra}{Anh}\ \ruby{Ri}{Lẫm}-\ruby{xư-ki}{Nguyệt}) là người gia nhập Nijisanji theo nhánh SEEDs, cùng với Banderas. Cô nàng là một công chúa tới từ hành tinh Hoa Anh Đào, với mục đích tới Trái Đất là để tìm lại bản thân.

Trái ngược với xuất thân nổi bật là một người ngoài hành tinh và công chúa tiên đến từ hành tinh khác, thực chất cô là một cô gái bình thường vụng về và năng động.

Cô nàng đđược đặc trưng bởi cách nói ngọng và giọng địa phương (sử dụng các cụm từ như ``\~{}ccha ne,'' ``\~{}ccha kedo,'' và ``\~{}yaken''). Lúc mới vào công ty, cô đã cố gắng kiềm chế giọng của mình hết mức có thể và dùng kính ngữ, nhưng sau một thời gian, cô nàng quyết định ``buông xõa'' và cứ nói như vậy bởi đó là tính cách của riêng mình, thậm chí còn sáng tạo thêm cả cách nói ``\~{}nbo'' vào cuối câu, nhưng dần dần cũng bỏ và thay bằng emoji hình cánh đào.

Ritsuki thường đạt được kết quả tốt trong các giải đấu và những giải tương tự, và mọi người thowfogn xem cô là ngừo "giỏi bất kỳ trò chơi nào", nhưng cô ấy tự nhận rằng cô không phải là người có khiếu thẩm mỹ về trò chơi, và tôi là kiểu người tiến bộ chậm, vì vậy tôi bù đắp bằng cách chơi nhiều. Xô ấy có tính cách nghiêm túc và thường luyện tập các trò chơi cạnh tranh ở một mức độ nào đó ở hậu trường trước khi chơi.

Cô nàng thích chơi chữ, tới độ còn\dots\ tổ chức một giải riêng về chơi chữ chỉ để nghe cách mọi người cười đùa với câu nói của mình thế nào, nhưng trớ trêu thay, Ritsuki lại\dots\ kém về khoản nghĩ ra cách để chơi chữ.

Ritsuki có mái tóc màu tím dài đến quá hông, được buộc hai bên thấp, trên mái có cài một cánh nơ màu đen ở bên trái và những cài tóc hình hoa anh đào ở xung quanh, đằng sau đầu có buộc một chiếc ruy băng trắng. Bộ trang phục cô nàng mặc là bộ truyền thống của Nhật Bản, gồm một bộ kimono trễ vai màu đen (hoặc hồng, tùy theo bộ cô mặc) mang họa tiết hoa anh đào, bên trong là một lớp váy trắng - tím bên trong được giữ bằng dây yếm hình chữ thập và một chiếc đai obi lớn màu vàng có ruy băng màu cam. Cô nàng đeo một đai buộc đùi ở bên trái và đi dép xăng-đan cao gót màu đen.

\chrint

\begin{center}
    \uline{Phần dưới đây là bài thơ, viết theo góc nhìn của Kaigou.}
\end{center}

\begin{center}
    Ngày thứ năm, sáng còn se,\\
    Ở nhà Viễn\footnote{Xem lại chú thích ở {\flaregothic ÉPISODE 126}, quyển số Mười - hồi Một.} bận phe phẩy chùi.\\
    Linh Mộc\footnote{Tức Suzuki, lấy từ 鈴木.} lo nồi nước sôi,\\
    Chẳng ai rảnh rỗi ngoài tôi rõ ràng.\\
    ``Cầu Vồng'' gọi, nhờ tay mang,\\
    Mấy tủ sách cũ dời sang kho gần.\\
    Bảo rằng thư viện bừa ngăn,\\
    Muốn tôi phụ giúp sắp dần từng khuôn.\\
    Tôi thì vốn ghét đàn ông,\\
    Không đời nào muốn dây dưa với họ.\\
    Nhưng Viễn ngó tôi nhè nhẹ:\\
    ``Cố giúp một hôm có sao đâu mà.''

    Tôi gật cho có lệ qua,\\
    Dặn điều kiện rõ: chỉ là nữ thôi.\\
    Không cho thằng nhãi nào ngồi,\\
    Không cho lão gác lại lôi thôi gần.\\
    Mắt ông Ban-đê-rát đăm,\\
    Chắc nghe đồn sẵn về tâm tính tôi.\\
    Ông ta chỉ gật, chép môi:\\
    ``Ừ, miễn đừng phá là tôi chẳng phiền.''\\
    Việc xong sớm, chẳng còn gì,\\
    Tôi vào góc khuất tìm ghế nghỉ chân.\\
    Vác mấy đống sách dày cuồn cuộn,\\
    Ngồi đọc miệt mài cho lùa cả ngày.

    \img{nj-5a}

    Tưởng rằng yên ổn một ngày,\\
    Ai dè phía cửa vang ngay tiếng người.\\
    Một thằng vạch đỏ phơ phơi,\\
    Mặt mũi rạng rỡ như chơi đại vờn.\\
    Lững thững bước, chẳng ai mời,\\
    Tựa vào giá sách như nơi quen rồi.\\
    Ghé tai cậu xoẹt xanh ngồi,\\
    Rỉ rả vài chữ, rồi cười nửa môi.\\
    Chẳng cần phải nói thêm lời,\\
    Tôi đoán được ngay: lại trò khua chiêng.\\
    Gì chứ đàn ông ngứa miệng,\\
    Chẳng mấy chốc sẽ khiến om sòm.

    \img{nj-5b}

    Ngồi bên giá sách im om,\\
    Có cậu hách-cơ cũng không khá gì.\\
    Giả vờ chăm chú mỗi khi,\\
    Nhưng chân cứ lắc như điệu kịch câm.\\
    Hai đứa rủ rỉ lầm rầm,\\
    Rồi bụm miệng cười, tôi nằm không yên.\\
    Chúng lén tháo gáy sách trên\\
    Gắn thêm mẩu giấy dán tên ngớ ngẩn.\\
    Một cuốn ghi ``Trà sữa cần'',\\
    Cuốn kia thì viết: ``Hôn nhân thất thường.''\\
    Tôi liếc qua, thấy mà thương,\\
    Mà cũng bực dọc, chẳng nhường ai đâu.

    Lại thêm cái giọng trêu nhau,\\
    Rì rầm giữa góc đằng sau tủ dài.\\
    Chúng đẩy sách như gài bẫy,\\
    Chỉ mong có đứa giở ra là tòe.\\
    Đến khi một cậu tóc hoe,\\
    Không biết từ đâu len về tận đây.\\
    Tay ôm cả xấp sách dày,\\
    Lại bị dính bẫy, ngã ngay đè lên.\\
    Tôi thở dài, chẳng muốn chen,\\
    Nhưng vẫn phải nhắc một phen cho rồi:\\
    ``Làm ơn giữ trật tự thôi,\\
    Ông già mà biết là tôi mặc kệ!''

    Na\footnote{Tức Akina, lấy từ chữ 那.} vỗ vai cậu Khôi\footnote{Tức Kai, lấy từ 灰.} ngồi,\\
    Khôi quay đầu lại, chưa kịp nói ra.\\
    Na liền bày tỏ thật thà\\
    Mặt mếu, môi nhếch như là thằng hề.\\
    Một tay nhéo mũi nhăn nhè,\\
    Một tay kéo tóc, mắt lòe như sương.\\
    Mặt như cái thước cong đường,\\
    Khiến cho tiếng phì bật vương quanh phòng.\\
    Na thì vẫn cố làm rồng,\\
    Méoméo chính mặt, hề không chịu thua.\\
    Khôi nhìn, nhịn cười chưa vừa,\\
    Tôi thì nén lại, mặt thưa thớt mày:\\
    ``Cái bản mặt ấy buồn thay!''\\
    Khúc khích khe khẽ như mây rớt về.\\
    Biết rằng chớ có đùa nghê,\\
    Ban mà nghe thấy là ê chề liền\dots

    Nguyệt\footnote{Tức Ritsuki, lấy từ chữ 月.} anh Đào liền dòm nghiêng,\\
    Nhìn ba người ấy đang rình rập nhau.\\
    Mặt Na đơ hệt miếng tàu,\\
    Tay còn cầm sách, vẫn nhàu nát tim.\\
    Không chỉ mặt méo ngả nghiêng,\\
    Na còn nhảy nhót, điên phiền mà vui.\\
    Tay thì chỉ thẳng lên trời,\\
    Múa điệu ma quái khiến người phát run.\\
    Khôi cùng tôi mặt lạnh trơn,\\
    Không dám thốt tiếng kẻo hơn phiền hà.\\
    Nguyệt kia chẳng chịu được mà\\
    Khúc khích bật tiếng, không ra thành lời!\\

    \img{nj-5c}

    Một luồng sét chớp sáng ngời,\\
    Lóe lên trong mắt khiến đời lặng im.\\
    Mồ hôi cô nhỏ thành tim,\\
    Giật mình hét khẽ, mặt chìm sắc tro:\\
    ``Dừng tay cả lũ ngốc to!\\
    Ban đang tiến lại, không lo đọc gì!?''\\
    Ban-đê-rát lập tức đi\\
    Dáng đi sừng sững, nghi nghi ngó nhìn.\\
    Đến nơi chẳng thấy tội tình,\\
    Chỉ toàn im lặng sách in lặng thầm.\\
    Mắt ai cũng hướng chuyên tâm,\\
    Không ai dám thở, chân nằm tại chỗ.\\
    Cả phòng như thể bị rồ,\\
    Chỉ mong thư quản thôi vô rồi về\dots

    Sóng gió tan, ai cũng thề,\\
    Cả bọn cùng thở dài lê thê lòng.\\
    Khôi quay nhắc bạn vừa xong:\\
    ``Lần sau bớt tếu, kẻo mong yên đời.''\\
    Tôi nhún nhẹ, nhoẻn vai cười,\\
    ``Nhưng mà nghĩ lại, buồn cười thật nha.''\\
    Nguyệt chen vào chào hỏi qua,\\
    Cũng hứa thầm sẽ tránh xa trò đùa.

    \img{nj-5d}

    Lặng yên một lúc chẳng thừa,\\
    Chợt nghe lạch cạch như mưa  hàng ngàn.\\
    Anh Đào bước nhẹ, mắt làn,\\
    Nheo nheo rồi khẽ ``ủa'' vang một lời.\\
    Thấy Khôi đang tỉ mẩn ngồi,\\
    Dùng kim, dùng kéo, ráp chơi chùa tròn.\\
    Na cùng tôi ngó lom khom,\\
    Trố tròn hai mắt, rồi còn bật cười.\\
    Tôi lùi lại một bước thôi,\\
    Nghĩ thầm: ``Không dính mớ trò kia đâu.''\\
    Nhưng ba cái đứa kia sau,\\
    Vẫn còn phá rối chốn lâu yên bình.

    \img{nj-5e}

    Tay Khôi vẫn đọc quyển kinh,\\
    Anh Đào dán mắt, chăm nhìn từng trang.\\
    Na thì hai tay giơ ngang,\\
    Gọi hai người bạn rời sang phía mình.\\
    Cậu xoay bút, động tác tinh,\\
    Một tay xoay nhẹ như hình cối xay.\\
    Nhắm tịt mắt, vẫn xoay hoài,\\
    Khiến cả ba đứa lặng người mà xem.\\
    Khôi thì cười nhẹ, mắt êm,\\
    Anh Đào thì vỗ tay thêm mấy lần.\\
    Tôi ngồi lặng ngắm vài phân,\\
    Mỉm cười khe khẽ: ``Tên này cũng liều!''

    \img{nj-5f}

    Na vừa diễn xong, Khôi liền\\
    Ngoắc tay bảo: ``Lại gần xem cái này!''\\
    Tôi còn chưa kịp gật ngay,\\
    Thì Khôi đã bước tới ngay cửa phòng.\\
    Cửa sổ mở rộng trông thông,\\
    Tòa tháp Đông Kinh đỏ hồng ngoài xa.\\
    Trời xanh vẫn sáng chan hòa,\\
    Khôi không nói, chỉ búng qua một lần.\\
    Chớp mắt, đèn tháp bừng phân,\\
    Cầu vồng hiện sáng muôn phần rực lên.

    \img{nj-5g}

    Anh Đào, An (Na) vỗ tay liền,\\
    Khen không ngớt miệng: ``Tuyệt nhiên hay mà!''\\
    Tôi nhìn, tòa tháp thật ra\\
    Cũng đành gật gù: ``Đẹp\dots\ nhưng thua Cường!''\\
    Na giật mình, lạnh sống lưng,\\
    Mắt tròn trợn ngược, tay bưng miệng cười.\\
    Lộp cộp tiếng bước chân rơi,\\
    Cả ba hấp tấp rối bời chạy ngay.\\
    Khôi thì lo lắng tỏ bày,\\
    ``Là Đào hay là Ban đây vậy trời?''\\
    Bỗng nghe tiếng thở hậm hừ,\\
    Ban-đê lại xuống, ngẩn ngơ đảo đầu.\\
    Chân ông sải bước thật mau,\\
    Quay về tiếp tân chẳng thể ngó ai.\\
    Ba người thở phảo nhẹ vai,\\
    ``Phen này may mắn, không ai bị trừ.''\\
    Tôi thì cười khẽ lắc đầu,\\
    ``Trò này đúng thật vui đâu kém gì.\\
    Nhưng mà đây chốn thư tri,\\
    Mà phát ra tiếng, là đi toi đời!\\
    Cấm phá rối, cấm chơi lơi,\\
    Lỡ ai bị đuổi, ai mời về thay?

    \img{nj-5g2}

    Tôi vừa dứt tiếng răn đay,\\
    Thì Đào đã nghịch chẳng nghe một lời.\\
    Mắt nàng lóe sáng tuyệt vời,\\
    Hai tay vẫy vẫy, miệng cười hí hô.\\
    Na cùng Khôi ngoảnh đầu dò,\\
    Ngờ đâu ``biến hình'' cô trò tấn công.\\
    Một vòng xoay, sáng bừng trong,\\
    Xanh xanh, hồng hồng, ánh lấp bềnh bồng.\\
    Cả thân phủ ánh kim đồng,\\
    Rầm rầm tiếng máy vang phòng rộn lên.\\
    Một tia chiếu thẳng từ trên,\\
    Sáu nòng súng lớn, đạn tên lửa dài.\\
    Cô nàng vũ trang chỉnh tề,\\
    Như vừa rút kiếm để về ra tay.\\
    Na và Khôi mắt ngất ngây,\\
    Vỗ tay rôm rốp ``Quá hay! Đỉnh rồi!''

    \img{nj-5h}

    Ngờ đâu chưa phải hết trò,\\
    Cả người Đào sáng rực hồng, bừng lên.\\
    La-de chớp, pháo hoa rền,\\
    Nổ vang như thể giữa đêm Giao thừa.\\
    Tôi ngồi lặng lẽ say sưa,\\
    ``\textit{Trò này đúng thật xứng vừa nhất năm!}''

    \begin{figure}[H]
        \centering
        \begin{subfigure}[t]{0.45\textwidth}
            \centering
            \includegraphics[width=8cm]{../../../../Image/Book?/nj-5i1.png}
        \end{subfigure}
        \quad
        \begin{subfigure}[t]{0.45\textwidth}
            \centering
            \includegraphics[width=8cm]{../../../../Image/Book?/nj-5i2.png}
        \end{subfigure}
    \end{figure}

    Bỗng vang lên tiếng vỗ tay,\\
    Không từ Khôi, cũng chẳng ai bày trò.\\
    Không từ Na, chẳng phải tôi,\\
    Cũng không phải Đào cười đòi khen thêm.\\
    Mọi người ngơ ngác đứng xem,\\
    Một người trông quầy bước đến âm thầm.\\
    Anh Đào mới cười bâng quơ,\\
    Chợt giật mình, mắt đổ dồn về Ban.\\
    Khôi thì trợn mắt ngỡ ngàng,\\
    Na run lẩy bẩy: ``Không làm gì đâu\dots''\\
    Nàng Đào lấm lét quay đầu,\\
    Thấy ông thư quản, râu rìa như rơm.\\
    Bàn tay đập bốp rõ to,\\
    Miệng cười đểu cáng, mắt loé như đèn.

    \img{nj-5j}

    ``Ban-đê-rát! Em xin thề\dots!''\\
    Đào run nói nhỏ, tìm lời giải vây.\\
    Nhưng ông chẳng để ai bày,\\
    Quát to: ``\textbf{Còn nói? Tới đây ăn đòn!}\\
    \textbf{Cái thư viện để học hành,\\
        Chứ đâu phải chỗ chị anh phá rầy!}''\\
    Khôi thì tái mặt như mây,\\
    Na xanh như lá úa ngày đầu đông.\\
    Còn tôi chẳng phải kẻ hòng \\
    Chỉ ngồi thở nhẹ, lặng trông ba người.\\
    ``Cái thư viện vốn rất đời,\\
    Chẳng phải cái chợ ai ơi biết điều!''\\
    Ban-đê-rát đập bàn kêu:\\
    ``\textbf{Chuẩn bị hối hận đi nhiều nhé con!}''\\
    Ông liền quát mắng dồn dồn,\\
    Dùng toàn ngôn ngữ ``trong son sắc màu''.\\
    Còn tôi rút nhanh về sau,\\
    Thở than: ``Không biết Viễn sao bây giờ\dots''\\

    \img{nj-5k}

    Chiều quay lại, thấy bất ngờ,\\
    Ba người hồi sáng bơ phờ ngục treo.\\
    Tay giơ, bị buộc dây theo,\\
    Cổ đeo bảng lỗi, nhìn lè lưỡi luôn.\\
    Đào thì bảng viết vuông tròn:\\
    ``Vì pháo hoa đẹp nên còn bị phê.''\\
    Na ghi: ``Tôi xúi giục kìa,\\
    Cho nên giờ mới ra nông nỗi này.''\\
    Khôi thì bảng đậm dòng ngay:\\
    ``Tôi là người dựng cái này - cầu vồng.''\\
    Thấy mà học hỏi đôi dòng:\\
    Thư viện, xin chớ làm không đúng chừng!\\
    Chốn tri thức, giữ cho cùng,\\
    Đừng đem trò cười quấy tung đất trời.\\
    Nếu ai cố ý phá chơi,\\
    Thì xem ba đứa rã rời chưa kia?\\

    \begin{figure}[H]
        \centering
        \includegraphics[width=12.5cm]{../../../../Image/Book?/nj-5l.png}
        \caption{Từ trái sang: Anh Lẫm Nguyệt (Sakura Ritsuki) với tấm bảng: ``Tôi là người khai màn pháo hoa đẹp đẽ'', Tam Chi Minh Na (Saegusa Akina) với tấm bảng: ``Tôi là người xúi giục'' và Đại Khôi (Mayuzumi Kai) với tấm bảng: ``Tôi tạo ra tháp RGB (đầy sắc màu)''.}
    \end{figure}

\end{center}
%==========================================TRUYỆN PHỤ=========================================
\omk

\begin{center}
    Tối về kể lại đôi lời,\\
    Cho Viễn với Mộc Linh ngồi nghe qua.\\
    Tôi vừa nói đến đoạn ba\\
    Tên kia bị bắt, Viễn ta thở dài.\\
    Cậu than: ``Trời ạ, đúng rồi\dots\\
    Y như hồi hô-lô chơi nhốn nhào.\\
    Mỗi lần định tập trung vào,\\
    Là y như rằng có bao chuyện phiền.\\
    Đứa thì hát, đứa điên điên,\\
    Đứa quậy tung cửa, lật nghiêng bàn làm.\\
    Tôi từng phát rồ không cam,\\
    Nhưng rồi cũng phải học làm thinh thôi\dots''

    Linh thì cười, mắt ngời ngời:\\
    ``Tội ba người đó quá trời là thương.\\
    Chắc giờ hối lỗi đầy gương,\\
    Nhưng mà cũng phải nhớ đường học luôn.''\\
    Viễn thì vẫn dựa sau lưng,\\
    Tay xoa trán nhẹ, mặt chừng chừng căm.\\
    Anh lẩm bẩm: ``Lần sau nhằm\\
    Đúng cái hôm nghỉ, đừng nhằm lúc tôi\dots''\\
    Tôi nghe thì cũng phì cười,\\
    Dù ban đầu cũng hơi hơi xấu hổ.\\
    Cái thư viện rõ là khổ,\\
    Có ba trò hề cũng khốn người ta\dots

    Một ngày nữa đã trôi qua,\\
    Vùng Cầu Vồng vẫn thẳm sâu chuyện người.\\
    Có khi lặng, có khi cười,\\
    Có khi phá phách cũng vui đôi phần.\\
    Tôi ngồi ghi lại từng vần,\\
    Viễn thì ngáp nhẹ, mắt dần lim dim.\\
    Suzu ngồi cuộn chăn im,\\
    Ngoài trời, sao rụng lặng im nỗi lòng.
\end{center}

\end{document}